\subsection{A common analytic comparison for circle shifts}

The word construction has supplied cheap translations and an expensive
half-turn. We now derive the inequality that compares those costs in an
arbitrary $\ell_1$ image. The proof is stated for any finite circle product;
the next subsection will verify its arithmetic hypothesis for the translation
subgroup $C_h$.

For a finite set $X$ and $A\subset X$, write
$\delta_A(x,y)=|\mathbf1_A(x)-\mathbf1_A(y)|$. The finite cut
representation of $\ell_1$ distances is standard
\cite[Section~4]{LinialLondonRabinovich1995}. We include the summability
argument because an $\ell_1$ target may have infinitely many coordinates.

\begin{lemma}[Finite cut representation]\label{pre:cuts}
For $f:X\to\ell_1$ with $X$ finite, there are nonnegative coefficients
$a_A$ on the nonempty proper subsets of $X$ such that
\[
 \|f(x)-f(y)\|_1=\sum_A a_A\delta_A(x,y),\qquad
 \sum_Aa_A\leq\sum_{\{x,y\}\subset X}\|f(x)-f(y)\|_1.
\]
\end{lemma}
\begin{proof}
For real $s,t$,
$|s-t|=\int_{\mathbb R}|\mathbf1_{s>u}-\mathbf1_{t>u}|\,du$.
Apply this to each coordinate of $f$ and group threshold sets that define
the same subset of $X$. Empty and full cuts contribute zero. Every other
cut separates some unordered pair, so its weight is counted at least once
in the displayed pair sum, which is finite. Nonnegative sums and integrals
may be interchanged, giving the representation and the bound.
\end{proof}

For this common lemma, let $M\geq2$ be any even integer,
$G=(M^{-1}\mathbb Z)/\mathbb Z$, and $V=G^r$ for
$r\geq1$. For a map $F:V\to\ell_1$ define its averaged displacement by
\[
 \Delta_F(a)=\E_{v\in V}\|F(v+a)-F(v)\|_1.
\]
All averages on finite groups in this paper are uniform. Represent a
character of $V$ as
$\chi_\alpha(v)=\exp(2\pi\mathrm i\sum_i\alpha_i v_i)$, choosing each
integer $\alpha_i$ with $|\alpha_i|\leq M/2$. Either representative at an
endpoint gives the same absolute value. Put
$T(\alpha)=\sum_i|\alpha_i|$ and
$\mathcal D_M=\{1,2,4,\ldots\}\cap[1,M/2]$.
For $\theta\in G$, write $\theta e_i$ for the vector in $V$ with value
$\theta$ in coordinate $i$ and zero elsewhere.

\begin{lemma}[Weighted finite-circle inequality]\label{pre:weighted-circle}
Let $H\leq V$, let $w=(1/2,\ldots,1/2)\in V$, and suppose $B>0$ satisfies
\begin{equation}\label{pre:branching-premise}
 \chi_\alpha|_H=1,\quad \chi_\alpha(w)\ne1
 \quad\Longrightarrow\quad T(\alpha)\geq B.
\end{equation}
For every $F:V\to\ell_1$ there are nonnegative weights $\omega_\alpha$ with
\begin{equation}\label{pre:fourier}
 \Delta_F(a)=\sum_\alpha\omega_\alpha|1-\chi_\alpha(a)|^2
 \qquad(a\in V).
\end{equation}
Writing $H^\perp=\{\alpha:\chi_\alpha|_H=1\}$, these weights satisfy
\begin{align}
 \sum_{\alpha\notin H^\perp}\omega_\alpha
   &=\tfrac12\E_{u\in H}\Delta_F(u),\label{pre:off-mass}\\
 4\sum_\alpha\omega_\alpha T(\alpha)
   &\leq\sum_{i=1}^r\sum_{a\in\mathcal D_M}
        \frac Ma\Delta_F((a/M)e_i).\label{pre:weighted-mass}
\end{align}
Consequently
\begin{equation}\label{pre:displacement}
 \Delta_F(w)\leq2\E_{u\in H}\Delta_F(u)
 +\frac1B\sum_{i=1}^r\sum_{a\in\mathcal D_M}
        \frac Ma\Delta_F((a/M)e_i).
\end{equation}
\end{lemma}
\begin{proof}
For a scalar function $g$ on $V$, normalized Fourier coefficients and
Parseval give
\[
 \E_v|g(v+a)-g(v)|^2
 =\sum_\alpha|\widehat g(\alpha)|^2|1-\chi_\alpha(a)|^2.
\]
For $g=\mathbf1_A$ its squared difference is the cut distance. Multiply
by the coefficients of Lemma~\ref{pre:cuts} and sum over cuts. This gives
\eqref{pre:fourier} with
$\omega_\alpha=\sum_Aa_A|\widehat{\mathbf1_A}(\alpha)|^2\geq0$.
The sums are finite after grouping cuts. Since a nontrivial character has
mean zero on a finite subgroup, the mean of $|1-\chi_\alpha(u)|^2$ on $H$
is two off $H^\perp$ and zero on it. This proves~\eqref{pre:off-mass}.

For each nonzero coordinate, take the largest power of two
$a\leq M/(2|\alpha_i|)$. It belongs to $\mathcal D_M$ and obeys
\[
 \frac14<\frac{a|\alpha_i|}{M}\leq\frac12,
 \qquad |1-e^{2\pi\mathrm i a\alpha_i/M}|^2\geq2,
 \qquad \frac Ma\geq2|\alpha_i|.
\]
That single scale contributes at least $4|\alpha_i|$. Sum over coordinates
and then against the nonnegative Fourier weights to obtain
\eqref{pre:weighted-mass}.

The half-turn multiplier is zero or four. Its detecting frequencies off
$H^\perp$ contribute at most $4\sum_{\alpha\notin H^\perp}\omega_\alpha$.
By~\eqref{pre:branching-premise}, the remaining ones contribute at most
$(4/B)\sum_\alpha\omega_\alpha T(\alpha)$. The two preceding bounds give
\eqref{pre:displacement}. Equivalently, the same estimates prove that
inequality for each character multiplier separately and then sum it with
nonnegative weights. Thus the lemma supplies both the weighted-mass and
the termwise-displacement forms of the comparison.
\end{proof}

The hypothesis~\eqref{pre:branching-premise} is the only arithmetic input
to this lemma. We shall prove it from the prime translations of each
construction. The inequalities themselves concern averages over all base
states; they impose no pointwise condition on individual Fourier weights.
